\section{Discussion}\label{sec:discussion}
The reduction keeps the temperature fixed. Computation is encoded in finite rational local tables, and the infinite-volume spin correlation detects whether a finite collision witness exists. Nonsurjectivity gives a strictly positive displacement from criticality. Surjectivity cancels every collision correction and leaves the same critical Ising marginal.

Theorem~\ref{thm:potts-main} adds a fixed order-parameter gap: the maximal magnetization is exactly $0$ or $m_9$, and the pressure jump is exactly $0$ or $9m_9/8$. Theorem~\ref{thm:main}, in contrast, has no total computable positive decay-rate bound, as Theorem~\ref{thm:modulus-main} makes precise. The Potts comparison uses the same witness, rewarded in both orientations; the rational critical odds are $3$. Nine is the smallest perfect square greater than four, giving both rational odds and a rigorously discontinuous transition.

All three constructions have the same fine-tuning limitation. The onto branch lies on a critical or coexistence surface of codimension one in the visible coupling parameters. Decreasing both Potts odds below $3$ makes every input subcritical, by the reward-free comparison. Thus arbitrarily small rational detuning can erase the dichotomy. The prescribed output-reset family preserves it because it changes only the collision retention and leaves that surface fixed. Uniform-replacement readout noise with retention $\eta>0$ multiplies $\mu(X_0)$ by $\eta$, and the distinct-site excess equal-colour correlation by $\eta^2$; this preserves the corresponding distinctions for the postprocessed observables.

Theorem~\ref{thm:full} gives full-state uniqueness and full clustering in the nonsurjective branch at $Q=257^2$. The free critical FK state has a witness-independent mass large enough to dominate two visible layers and the rare hidden selectors. In the surjective branch a disordered state with the same bound clusters, while ordered phases survive. Thus the new undecidable content is the number of phases, not a decay rate. The nine-colour and Ising uniqueness problems remain open.

\begin{remark}[Conjecture: first-order versus continuous transitions]
The contrast suggests a conjectural organizing principle: near a first-order point, undecidability may be qualitative, with uniform disordered-phase rates and one undecidable bit recording survival of ordered phases; near a continuous point it may be quantitative, carried by the lack of a computable decay modulus. The first statement is realized by Theorem~\ref{thm:full}, and the second by the Ising Theorem~\ref{thm:modulus-main}. A general correspondence between transition order and the form of undecidability is a conjecture, not a consequence of these two examples. In particular the proof supplies no uniform computable finite-volume boundary crossover scale.
\end{remark}

The all-state argument is separate from the free-hidden finite-volume expansion. Uniform conditional uniqueness of the replicas provides local insertion ratios in arbitrary infinite bond configurations. Wired-core comparison then controls every Gibbs state without imposing full-state uniqueness in the nonsurjective branch.

For the Ising and nine-colour families the difficult remaining issue is the response of even and hidden observables. The counterexamples show why short magnetic clusters and conditional hidden contraction alone are insufficient. A successful completion must use the collision identities and their quantitative smallness.

\begin{remark}[Noetherian obstruction for a fixed collision template]\label{rem:noether}
Fix a finite template with finitely many rational bias parameters $z=(z_1,\ldots,z_b)$, restricted to an explicitly decidable rational domain where all denominators below are nonzero. Suppose a branch is characterized by exact identities $R_j(z)=0$, $j\ge1$, with each $R_j$ a rational function over $\mathbb Q$. Clearing denominators gives polynomials $p_j$. The Hilbert basis theorem \cite[Chapter 2]{CLO} makes the ideal $(p_1,p_2,\ldots)\subset\mathbb Q[z]$ finitely generated. Moreover finitely many of the original $p_j$ suffice: each of finitely many ideal generators is a finite polynomial combination of original members; collect those members. Consequently all identities hold exactly when this finite sublist vanishes. Exact rational (or algebraic) evaluation decides the branch. This is an existence-of-a-decider statement for each fixed template; it supplies no uniform procedure for finding that finite sublist from an infinite presentation.

For one fixed CA $F$, with iid input biases in finitely many parameters, every $c_T$ is a rational polynomial in those biases: use the direct finite sum of integer-valued collision products. Apply the preceding argument to all nonempty $T$. A finite sublist fits in some box $[-N(F),N(F)]^d$ (the bias parametrization is fixed with $F$). Hence collision-nullity on all $T$ is equivalent to nullity for all nonempty subsets of that finite box. No fixed CA with rational biases can encode an undecidable predicate solely in this exact null locus.

With fixed alphabet, dimension and range there are only finitely many deterministic local rules, so the same obstruction covers that finite collection of collision templates. In a collision-nullity route with these restrictions, range must therefore grow with the input, as it does here. This conclusion concerns fixed finite templates and exact identities; it is not an obstruction to limit predicates, infinite inequalities, varying alphabets, or every possible fixed-range construction.
\end{remark}

\begin{remark}[Uniform fields and raw decay; robustness sketches]\label{rem:robust}
For translation-invariant positive finite-range Ising-spin models, raw all-state spin decay has empty interior in any real parameter space allowing a uniform field $H\sum_x\sigma_x$. Here is a short proof. A translation-invariant Gibbs state with summable raw spin correlations has zero magnetization: its squared spatial-average mean is bounded by the second moment of that average, which tends to zero by summability. On any compact interval of fields, finite energy bounds single-spin conditional variances below by a constant $c>0$. Choose a sublattice of positive density whose distinct sites occur in no common interaction term. Conditional on its complement these spins are independent. For $M_\Lambda=\sum_{x\in\Lambda}\sigma_x$, total variance gives
$\operatorname{Var}(M_\Lambda)\ge c'|\Lambda|$ up to boundary terms. Therefore the finite-volume pressures have second derivative at least $c'>0$, and their limit is strongly convex in $H$. Magnetizations of translation-invariant Gibbs states are its subgradients (the Gibbs variational principle \cite[Chapter 15]{Georgii}). Strong convexity allows subgradient zero at at most one field. Since raw all-state decay forces that zero subgradient, it can hold at at most one point of each uniform-field line, proving empty interior. Fields correspond to arbitrary positive real weights; the interior assertion uses this real topology.

For contrast, the following robust-region statements are sketches based on standard mixing and low-temperature phase theory \cite[Chapters 6--7]{FV}. Dobrushin's strict contraction condition gives an open region of exponential decay for connected correlations. It also gives raw decay for $D_x=\sigma_x-\sigma_{x+e_1}$, since uniqueness and translation invariance imply $\mu(D_x)=0$. At low temperature the square-lattice checkerboard antiferromagnet has two translation-related phases throughout an open interval of small uniform fields. Their symmetric mixture has persistent staggered correlations and nonzero long-distance $D$ correlations along even displacements, giving open regions of failure for the all-state connected and $D$-raw properties. Extension to small translation-invariant finite-range perturbations uses the Peierls/contour hypotheses of Pirogov--Sinai theory; that contour argument is not proved here. These sketches make no robustness claim about Theorem~\ref{thm:main}.
\end{remark}

\begin{remark}[Why computable-real weights make full clustering easy to encode]\label{rem:real}
Enumerate finite Garden-of-Eden witnesses for a CA $F$ by a fixed effective search, at integer times $T\ge0$, and set
\[
 \delta_F=\begin{cases}2^{-(T+2)},&\text{the first witness appears at time }T,\\
                       0,&\text{no witness appears}.\end{cases}
 \qquad e^{2K_H}=3-\delta_F,\quad e^{2K_V}=2.
\]
This is a uniformly computable real input: simulate through time $n$; if no witness has appeared, output zero with error at most $2^{-(n+3)}$. If one has appeared, its dyadic value is exact. Logarithms of these positive computable reals also give computable energy entries.

A witness exists exactly for nonsurjective $F$: a nononto continuous image of the compact full shift is closed and omits a cylinder; testing whether a finite pattern has a preimage is a finite local enumeration. If $\delta_F>0$, the odds product is $2-\delta_F<2$, so this nearest-neighbor ferromagnet is strictly subcritical and has a unique Gibbs state with exponential clustering of all bounded local observables. To recall the full-clustering step, (I3) and the Edwards--Sokal boundary coupling bound influence on an interior set by the probability of an FK connection to the boundary; exponential connectivity makes this influence exponentially small. Conditioning a local observable on a distant exterior and integrating then gives exponential covariance bounds for arbitrary finite supports \cite[Chapters 3 and 6]{FV}. If $\delta_F=0$, the fixed critical spin correlation already fails exponential decay. Thus deciding full clustering for these computable-real descriptions would decide CA surjectivity.

The numerical value $3-\delta_F$ is rational in each case, but its presentation here is a computable-real name: no algorithm outputs its exact rational numerator and denominator without deciding whether a witness will ever appear. The content of Theorem~\ref{thm:main} is its finite \emph{rational weight-table} input, not merely that the values happen to be rational.
\end{remark}
